\documentclass{article}
\usepackage{amsmath,amssymb}
\usepackage[vocab={numbers,arithmetic}]{leanify}
\newtheorem{lemma}{Lemma}
\begin{document}

\section{Notation}

Let $\varphi$ denote Euler's totient function.
We write $\mu$ for the M\"obius function.
We write $\tau(n)$ for the number of divisors of $n$.
We write $\sigma(n)$ for the sum of the divisors of $n$.
We write $\pi(x)$ for the number of primes not exceeding $x$.
We write $p_n$ for the $n$-th prime.
We write $\omega(n)$ for the number of distinct prime factors of $n$.
We write $\nu_p(n)$ for the $p$-adic valuation of $n$.
We write $\operatorname{ord}_n(a)$ for the multiplicative order of $a$ modulo $n$.
We write $s(n)$ for the sum of the digits of $n$.

\section{Arithmetic functions}

\begin{lemma}\label{totient}
For every prime $p$, $\varphi(p)=p-1$. Moreover, $\varphi(12)=4$ and $\mu(6)=1$ and $\mu(12)=0$.
\end{lemma}

\begin{lemma}\label{divisors}
$\tau(12)=6$ and $\sigma(6)=12$, and $6$ is perfect and $12$ is abundant.
\end{lemma}

\begin{lemma}\label{primes}
$\pi(10)=4$ and $p_5=11$ and $\omega(60)=3$ and $\nu_2(48)=4$.
\end{lemma}

\begin{lemma}\label{order}
$\operatorname{ord}_7(2)=3$ and $s(1234)=10$.
\end{lemma}

\begin{lemma}\label{multiplicative}
$\varphi$ is multiplicative and $\mu$ is multiplicative. Moreover, $\tau(n)\le2\sqrt{n}$ for every positive integer $n$.
\end{lemma}

\begin{lemma}\label{mersenne}
For every Mersenne prime $p$, $p\,\frac{p+1}{2}$ is a perfect number.
\end{lemma}

\begin{lemma}\label{prime-factorization}
Every integer $n\ge2$ is a product of primes.
\end{lemma}

\begin{lemma}\label{goldbach-ten}
$10$ is a sum of two primes and $7$ is a difference of two squares.
\end{lemma}

\begin{lemma}\label{lagrange-seven}
$7$ is a sum of four squares.
\end{lemma}

\begin{lemma}\label{not-both-zero}
Let $a$ and $b$ be integers, not both zero. Then $a^2+b^2>0$.
\end{lemma}

\begin{lemma}\label{pair-nonzero}
Let $a,b$ be integers with $(a,b)\neq(0,0)$. Then $a\neq0$ or $b\neq0$.
\end{lemma}

\begin{lemma}\label{non-zero}
Let $n$ be a non-zero integer. Then $n^2\ge1$.
\end{lemma}

\begin{lemma}\label{gcd-lcm}
For all positive integers $a$ and $b$, $\gcd(a,b)\cdot\operatorname{lcm}(a,b)=ab$.
\end{lemma}

\begin{lemma}\label{post-if}
Let $n\ge2$ be an integer. Then $n^2>n$ if $n$ is odd.
\end{lemma}

\begin{lemma}\label{every-has}
Every integer $n>1$ has a prime divisor.
\end{lemma}

\begin{lemma}\label{post-for}
Let $p$ be a prime. Then $p$ divides $k^p-k$ for $k>0$.
\end{lemma}

\begin{lemma}\label{written-as}
Every positive integer $n\ge2$ can be written as a product of primes.
\end{lemma}

\begin{lemma}\label{either-or}
Let $n$ be an integer. Then either $n$ is even or $n$ is odd.
\end{lemma}

\begin{lemma}\label{has-or}
Let $n\ge2$ be an integer. Then $n$ has a prime factor $p\le\sqrt{n}$ or $n$ is prime.
\end{lemma}

\begin{lemma}\label{pair}
Let $(a,b)$ be a pair of integers with $a<b$. Then $a+1\le b$.
\end{lemma}

\begin{lemma}\label{comparatives-and-sums}
Every integer greater than $1$ has a prime divisor. Every prime number greater than $2$ is odd. Let $n$ be an integer greater than or equal to $2$. Then $n^2>n$. The square of an odd integer is odd. The sum of two odd integers is even. The difference of two even integers is even.
\end{lemma}

\begin{lemma}\label{postposed-if}
Let $n$ be an integer. Then $n^2$ is even if $n$ is even.
\end{lemma}

\end{document}
